\section{Программирование светодиодов}
\subsection{Введение}
Этот урок вводит в программирование индикации на квадрокоптере «Геоскан Пионер» на языке Lua. Мы начнём с определений, затем — теория, примеры и практические задания.

\subsection{Толковый словарь}
Нажмите на термин в тексте, чтобы перейти к его определению.
\begin{description}
  \item[\hypertarget{term:module}{Модуль}] часть программы, отвечающая за конкретную задачу и предоставляющая функции для работы. Модуль \hyperlink{term:ledbar}{Ledbar} управляет светодиодами.
  \item[\hypertarget{term:driver}{Драйвер}] программный компонент, который общается с аппаратурой и переводит команды в сигналы для устройств.
  \item[\hypertarget{term:buffer}{Буфер}] область памяти, куда мы записываем значения (например, цвета); драйвер считывает буфер и выводит свечение.
  \item[\hypertarget{term:variable}{Переменная}] именованная «коробочка» со значением. Пример: \verb|x = 10|.
  \item[\hypertarget{term:scope}{Локальная/Глобальная}] локальная переменная объявляется через \verb|local| и видна только внутри текущего блока; глобальная (без \verb|local|) видна во всём скрипте.
  \item[\hypertarget{term:timer}{Таймер}] «будильник», который вызывает указанную функцию периодически или один раз с задержкой.
  \item[\hypertarget{term:period}{Период}] время между двумя вызовами таймера.
  \item[\hypertarget{term:tick}{Тик}] один вызов функции таймером.
  \item[\hypertarget{term:state}{Состояние}] текущее положение логики (например, «красный», «жёлтый», «зелёный»).
  \item[\hypertarget{term:automaton}{Конечный автомат}] модель, где программа находится в одном из состояний и переходит в другое при событиях.
  \item[\hypertarget{term:rgb}{RGB-компоненты}] три числа \verb|R|, \verb|G|, \verb|B|, задающие яркости каналов; каждое в \verb|[0.0, 1.0]|.
  \item[\hypertarget{term:index}{Индекс светодиода}] порядковый номер «лампочки» в линейке; используем \verb|0..24|.
  \item[\hypertarget{term:api}{API}] набор функций модуля (\verb|Ledbar.set|, \verb|Timer.new| и др.).
  \item[\hypertarget{term:ledbar}{Ledbar}] модуль управления RGB-светодиодами квадрокоптера.
\end{description}

\subsection{Аппаратная часть и модуль Ledbar}
\subsection{Физическая реализация}
На плате установлено \textbf{4 RGB-светодиода}. Каждый имеет три кристалла: красный (R), зелёный (G), синий (B). Комбинация трёх яркостей даёт итоговый цвет.

\subsection{Особенности адресации (Буфер драйвера)}
Важно: \hyperlink{term:ledbar}{Ledbar} работает с \hyperlink{term:buffer}{буфером} \textbf{из 25 элементов}. Поэтому инициализация должна задавать размерность \verb|25|.
\begin{quote}
\textbf{Требование:} \verb|Ledbar.new(25)|. Значение \verb|4| использовать нельзя — диоды не загорятся.
\end{quote}
\begin{minted}{lua}
local unpack = table.unpack
local ledNumber = 25                  -- размерность буфера
local leds = Ledbar.new(ledNumber)    -- объект управления лентой

local function changeColor(col)
  for i=0, ledNumber-1 do             -- индексы 0..24
    leds:set(i, unpack(col))          -- записываем RGB
  end
end

changeColor({0,1,0})                  -- включаем зелёный
\end{minted}

\subsection{Теория и интерфейсы API}
\subsection{RGB-модель и диапазоны}
\hyperlink{term:rgb}{RGB} — три управляющих числа: \verb|R| (красный), \verb|G| (зелёный), \verb|B| (синий) в \verb|[0.0, 1.0]|. \verb|0.0| — канал выключен, \verb|1.0| — горит ярко. Смешивая каналы, получаем цвета: \verb|{1,0,0}| (красный), \verb|{0,1,0}| (зелёный), \verb|{0,0,1}| (синий), \verb|{1,1,0}| (жёлтый), \verb|{1,1,1}| (белый), \verb|{0,0,0}| (выключено). Плавная смена значений позволяет создавать анимации.

\subsection{Адресация и буфер}
\hyperlink{term:index}{Индекс} — позиция \verb|0..24|. \hyperlink{term:buffer}{Буфер} — массив из 25 ячеек, где каждая хранит цвет своего индекса. Драйвер считывает весь буфер, поэтому обновляйте все элементы, даже если видны только первые четыре.

\subsection{Модуль Ledbar}
\begin{itemize}
  \item \verb|Ledbar.new(N)| — создаёт объект управления (\hyperlink{term:api}{API}); используйте \verb|25|.
  \item \verb|leds:set(i, r, g, b)| — задаёт цвет индекса \verb|i| (\verb|r,g,b| в \verb|[0.0,1.0]|).
  \item Вспомогательная \verb|changeColor(col)| перезаписывает все индексы.
\end{itemize}
\textbf{Размерность} — число элементов буфера (\verb|25|). \textbf{Индекс} — порядковый номер, начиная с \verb|0|.

\subsection{Таймеры}
\begin{itemize}
  \item \verb|Timer.new(period, fn)| — создаёт \hyperlink{term:timer}{периодический таймер}; \verb|period| — \hyperlink{term:period}{период} в секундах; \verb|fn| вызывается каждый \hyperlink{term:tick}{тик}.
  \item \verb|timer:start()|, \verb|timer:stop()| — запуск/остановка таймера.
  \item \verb|Timer.callLater(delay, fn)| — единичный вызов \verb|fn| через \verb|delay| секунд.
  \item Таймеры и \hyperlink{term:state}{состояния} делайте \hyperlink{term:scope}{глобальными} (без \verb|local|), иначе сборщик мусора их удалит.
\end{itemize}
\textbf{Пример периодического таймера (мигание):}
\begin{minted}{lua}
local ledNumber = 25
local leds = Ledbar.new(ledNumber)
local on = false
function updateBlink()
  local r,g,b = (on and 1 or 0), 0, 0
  for i=0, ledNumber-1 do leds:set(i, r, g, b) end
  on = not on
end
blinkTimer = Timer.new(0.5, updateBlink)  -- 0.5 секунды
blinkTimer:start()
\end{minted}
\textbf{Пример отложенного вызова (через 3 секунды):}
\begin{minted}{lua}
local ledNumber = 25
local leds = Ledbar.new(ledNumber)
Timer.callLater(3, function ()
  for i=0, ledNumber-1 do leds:set(i, 0, 0, 1) end  -- включить синий
end)
\end{minted}

\subsection{Событийная модель}
\begin{itemize}
  \item \verb|function callback(event)| — обработчик системных событий.
  \item Примеры: \verb|Ev.COPTER_LANDED| (посадка), \verb|Ev.LOW_VOLTAGE2| (низкое напряжение), \verb|Ev.POINT_REACHED|, \verb|Ev.ALTITUDE_REACHED|.
  \item Внутри \verb|callback| быстро реагируйте на событие; не используйте \verb|sleep()|.
\end{itemize}
\textbf{Примеры:}
\begin{minted}{lua}
local ledNumber = 25
local leds = Ledbar.new(ledNumber)
function callback(event)
  if (event == Ev.COPTER_LANDED) then
    for i=0, ledNumber-1 do leds:set(i, 0, 0, 0) end
  end
end
\end{minted}
\begin{minted}{lua}
local ledNumber = 25
local leds = Ledbar.new(ledNumber)
mainTimer = Timer.new(0.1, function () end)
mainTimer:start()
function callback(event)
  if (event == Ev.LOW_VOLTAGE2) then
    mainTimer:stop()
    Timer.callLater(1, function ()
      for i=0, ledNumber-1 do leds:set(i, 1, 0, 0) end
    end)
  end
end
\end{minted}

\subsection{Анализ подходов к программированию}
\subsection{Конечный автомат}
Подход основан на \hyperlink{term:automaton}{конечном автомате}: программа находится в одном из заранее определённых состояний, а переходы между ними происходят при наступлении системных событий или по внутренней логике.
Состояние хранится в переменной (например, \verb|curr_state|), а набор действий для каждого состояния определяется таблицей функций \verb|action|. 
При поступлении события (\verb|callback(event)|) выбирается текущая функция из \verb|action[curr_state]| и выполняется.
Этот подход удобен для сценариев «включить один раз и завершить», либо для дискретных, чётко структурированных шагов. 
См. пример структуры в файле \verb|4_светодиода.lua|: инициализация \verb|Ledbar.new(25)|, установка фиксированных цветов \verb|set(i, r, g, b)| и завершение работы.
\subsection{Блокирующее выполнение}
Блокирующий стиль использует задержки \verb|sleep()|, при которых выполнение программы останавливается. 
На борту квадрокоптера такой подход нежелателен: блокировка может мешать обработке других задач (навигация, связь, датчики). 
Вместо \verb|sleep()| следует применять \hyperlink{term:timer}{асинхронные таймеры}, которые вызывают функцию через заданные интервалы, не «замораживая» систему. 
См. файл \verb|Светодиоды_и_таймер.lua| как пример, где последовательные \verb|sleep| приводят к пошаговым включениям, но архитектурно безопаснее заменить их на \verb|Timer.new| или \verb|Timer.callLater|.
\subsection{Асинхронные таймеры}
Рекомендуемый подход: все анимации строятся на \verb|Timer.new(period, fn)|. 
Функция \verb|fn| выполняет «малый шаг» анимации (обновление цвета, переключение состояния), а таймер срабатывает регулярно с заданным периодом. 
Важно объявлять таймеры глобально, чтобы их не удалил сборщик мусора, и избегать длительных операций внутри \verb|fn|. 
См. \verb|nice_color (1).lua| как образец: глобальный таймер, небольшие приращения компонент цвета, полная перерисовка буфера \verb|0..24| на каждом тике.

\subsection{Примеры кода}
\begin{minted}{lua}
-- создание порта управления светодиодом
local ledbar = Ledbar.new(25)

-- переменная текущего состояния
local curr_state = "START"

-- таблица функций, вызываемых в зависимости от состояния
action = {
	["START"] = function(x)

		ledbar:set(0, 1, 0.0, 0.0)

		ledbar:set(1, 0.0, 1, 0.0)

		ledbar:set(2, 0.0, 0.0, 1)

		ledbar:set(3, 1, 1, 0.0)

		-- выключение двигателей и конец программы
		ap.push(Ev.ENGINES_DISARM)
		curr_state = "NONE"

	end,
}

-- функция обработки событий, автоматически вызывается автопилотом
function callback(event)
	if (event == Ev.ALTITUDE_REACHED) then
		action[curr_state]()
	end

	if (event == Ev.POINT_REACHED) then
		action[curr_state]()
	end

	if (event == Ev.COPTER_LANDED) then
		sleep(2)
		action[curr_state]()
	end
end

-- вызов функции из таблицы состояний, соответствующей первому состоянию
action[curr_state]()
\end{minted}

\begin{minted}{lua}
-- создание порта управления светодиодом
local ledbar = Ledbar.new(25)

-- переменная текущего состояния
local curr_state = "START"

-- таблица функций, вызываемых в зависимости от состояния
action = {
	["START"] = function(x)

		ledbar:set(0, 1, 0.0, 0.0)

		sleep(1000 / 1000.0) 
		ledbar:set(1, 0.0, 1, 0.0)

		sleep(1000 / 1000.0) 
		ledbar:set(2, 0.0, 0.0, 1)

		sleep(1000 / 1000.0) 
		ledbar:set(3, 1, 1, 0.0)

		-- выключение двигателей и конец программы
		ap.push(Ev.ENGINES_DISARM)
		curr_state = "NONE"

	end,
}

-- функция обработки событий, автоматически вызывается автопилотом
function callback(event)
	if (event == Ev.ALTITUDE_REACHED) then
		action[curr_state]()
	end

	if (event == Ev.POINT_REACHED) then
		action[curr_state]()
	end

	if (event == Ev.COPTER_LANDED) then
		sleep(2)
		action[curr_state]()
	end
end

-- вызов функции из таблицы состояний, соответствующей первому состоянию
action[curr_state]()
\end{minted}

\begin{minted}{lua}
-- https://learnxinyminutes.com/docs/ru-ru/lua-ru/ ссылка для быстрого ознакомления с основами языка LUA

-- Скрипт реализует мигание светодиодов на базовой плате случайными цветами

-- Упрощение вызова функции распаковки таблиц из модуля table
local unpack = table.unpack
-- Количество светодиодов на базовой плате
local ledNumber = 25
-- Создание порта управления светодиодами
local leds = Ledbar.new(ledNumber)

-- Функция, изменяющая цвет 4-х RGB светодиодов на базовой плате
local function changeColor(col)
    for i=0, ledNumber - 1, 1 do
        leds:set(i, unpack(col))
    end
end

-- Функция, реализующая выключение таймера и изменение цвета светодиодов на красный
local function emergency()
    -- Остановка таймера
    timerRandomLED:stop()
    -- Изменение цвета светодиодов (через секунду - т.к. после остановки таймера timerRandomLED
    -- его функция выполнится еще раз)
    Timer.callLater(1, function () changeColor({1, 0, 0}) end)
end

-- Функция обработки событий, автоматически вызывается автопилотом
function callback(event)
    -- Вызов функции emergency() при низком напряжении на аккумуляторе
    if (event == Ev.LOW_VOLTAGE2) then
        emergency()
    end
end

color = {0, 0, 0} 
step = 0.05 
index = 1

-- Создание таймера, каждую секунду меняющего цвета каждого из 4-х светодиодов на случайные
timerRandomLED = Timer.new(0.05, function ()
    color[index] = color[index] + step
    if color[index] > 0.95 or color[index] < 0.05 then
        if index < 3 then
            index = index + 1
        else
            index = 1
            step = -step
        end
    end
    changeColor(color)
end)
-- Запуск созданного таймера
timerRandomLED:start()
\end{minted}
\subsection{Проверка знаний}
Перед выполнением практических заданий ответьте на вопросы. Вставьте пропущенную строчку кода или слово. Выберите один верный вариант и сверьте с пояснением.
\begin{enumerate}
  \item Какое значение размера буфера должно быть передано в инициализацию? \verb|leds = Ledbar.new(___)|
  \begin{enumerate}[label=(\alph*)]
    \item 4 \item 8 \item 25 \item 1
  \end{enumerate}
  Ответ: (c). Пояснение: драйвер считывает буфер из 25 элементов; значение 4 не работает.

  \item Чем должно быть начальное значение счётчика в цикле? \verb|for i = ____, ledNumber-1 do ... end|
  \begin{enumerate}[label=(\alph*)]
    \item 0 \item 1 \item -1 \item ledNumber
  \end{enumerate}
  Ответ: (a). Пояснение: индексы линейки начинаются с 0 и заканчиваются 24.

  \item Каким вызовом передать RGB-таблицу в \verb|leds:set| при объявленном \verb|local unpack = table.unpack|?
  \begin{enumerate}[label=(\alph*)]
    \item \verb|leds:set(i, col)| \item \verb|leds:set(i, unpack(col))| \item \verb|leds:set(i, toRGB(col))| \item \verb|leds:set(i, rgb(col))|
  \end{enumerate}
  Ответ: (b). Пояснение: \verb|unpack| распаковывает таблицу \verb|{r,g,b}| в три аргумента.

  \item Как корректно единожды включить зелёный цвет на всей линейке?
  \begin{enumerate}[label=(\alph*)]
    \item \verb|changeColor({0,1,0})| \item \verb|changeColor({1,0,0})| \item \verb|changeColor({0,0,1})| \item \verb|changeColor({1,1,1})|
  \end{enumerate}
  Ответ: (a). Пояснение: компоненты RGB для зелёного — \verb|{0,1,0}|.

  \item Как создать отложенный запуск синего через 3 секунды?
  \begin{enumerate}[label=(\alph*)]
    \item \verb|sleep(3); changeColor({0,0,1})| \item \verb|Timer.new(3, changeColor)| \item \verb|Timer.callLater(3, function() changeColor({0,0,1}) end)| \item \verb|ap.push(Ev.COPTER_LANDED)|
  \end{enumerate}
  Ответ: (c). Пояснение: \verb|Timer.callLater| планирует единичный вызов без блокировки.

  \item Какой период задать для мигания раз в полсекунды? \verb|Timer.new(___, fn)|
  \begin{enumerate}[label=(\alph*)]
    \item 0.5 \item 500 \item 50 \item 1.5
  \end{enumerate}
  Ответ: (a). Пояснение: период указывается в секундах.

  \item Как инициализировать булевое состояние для мигания?
  \begin{enumerate}[label=(\alph*)]
    \item \verb|on = false| \item \verb|on = 0| \item \verb|on = "false"| \item \verb|on = nil|
  \end{enumerate}
  Ответ: (a). Пояснение: булев тип Lua использует \verb|true/false|.

  \item Как проверить чётность индекса \verb|i|?
  \begin{enumerate}[label=(\alph*)]
    \item \verb|i / 2 == 0| \item \verb|i % 2 == 0| \item \verb|i == even| \item \verb|isEven(i)|
  \end{enumerate}
  Ответ: (b). Пояснение: оператор остатка \verb|%| определяет чётность.

  \item Как остановить запущенный таймер \verb|t|?
  \begin{enumerate}[label=(\alph*)]
    \item \verb|t:stop()| \item \verb|t:pause()| \item \verb|stop(t)| \item \verb|t.cancel()|
  \end{enumerate}
  Ответ: (a). Пояснение: API таймера предоставляет методы \verb|start/stop|.

  \item Как запустить таймер \verb|t| после создания?
  \begin{enumerate}[label=(\alph*)]
    \item \verb|t:run()| \item \verb|t:start()| \item \verb|start(t)| \item \verb|t:on()|
  \end{enumerate}
  Ответ: (b). Пояснение: корректный метод — \verb|start|.

  \item Как называется функция-обработчик событий автопилота?
  \begin{enumerate}[label=(\alph*)]
    \item \verb|function events(e)| \item \verb|function callback(event)| \item \verb|function onEvent(e)| \item \verb|function handler(ev)|
  \end{enumerate}
  Ответ: (b). Пояснение: платформа вызывает \verb|callback| и передаёт \verb|event|.

  \item Какое событие соответствует посадке?
  \begin{enumerate}[label=(\alph*)]
    \item \verb|Ev.POINT_REACHED| \item \verb|Ev.ALTITUDE_REACHED| \item \verb|Ev.COPTER_LANDED| \item \verb|Ev.LOW_VOLTAGE2|
  \end{enumerate}
  Ответ: (c). Пояснение: \verb|COPTER_LANDED| — системное событие посадки.

  \item Какое событие соответствует низкому напряжению?
  \begin{enumerate}[label=(\alph*)]
    \item \verb|Ev.LOW_VOLTAGE2| \item \verb|Ev.BATTERY_OK| \item \verb|Ev.POWER_ON| \item \verb|Ev.MOTOR_ARMED|
  \end{enumerate}
  Ответ: (a). Пояснение: используйте \verb|LOW_VOLTAGE2| для аварийной индикации.

  \item Как выглядит условие для «полосы загрузки»: включить все индексы \verb|<= k|?
  \begin{enumerate}[label=(\alph*)]
    \item \verb|if i < k then ...| \item \verb|if i <= k then ...| \item \verb|if k <= i then ...| \item \verb|if i == k then ...|
  \end{enumerate}
  Ответ: (b). Пояснение: включёнными остаются все индексы не выше текущего \verb|k|.

  \item Какой допустимый диапазон значений \verb|r,g,b| в модели RGB?
  \begin{enumerate}[label=(\alph*)]
    \item \verb|0..255| \item \verb|0..100| \item \verb|0.0..1.0| \item \verb|-1..1|
  \end{enumerate}
  Ответ: (c). Пояснение: платформа принимает вещественные значения от \verb|0.0| до \verb|1.0|.
\end{enumerate}

\subsection{Задания для лабораторной работы}
Используйте асинхронные таймеры, глобальные переменные для таймеров и состояний, полный проход по индексам \verb|0..24|. Формулировки ниже предельно простые; рядом указаны подсказки.
\subsection{Базовые алгоритмы}

\begin{enumerate}
  \item \textbf{Одновременное включение всей светодиодной линейки зелёным цветом.}  
  Требуется реализовать алгоритм, обеспечивающий однократную установку зелёного цвета с компонентами \verb|{0,1,0}| на всех светодиодах линейки с индексами от \verb|0| до \verb|24|. Для этого необходимо создать объект светодиодной линейки, реализовать функцию массовой установки цвета и выполнить её вызов.  
  \textit{Ожидаемый результат:} вся линейка равномерно светится зелёным цветом.

  \item \textbf{Отложенное включение синего цвета.}  
  Необходимо реализовать алгоритм, при котором светодиодная линейка остаётся выключенной в течение первых трёх секунд работы программы, после чего все светодиоды одновременно загораются синим цветом \verb|{0,0,1}|. Реализация должна использовать механизм отложенного вызова без применения блокирующих задержек.  
  \textit{Ожидаемый результат:} через 3 секунды после запуска программы линейка полностью загорается синим цветом.

  \item \textbf{Периодическое мигание красным цветом.}  
  Требуется реализовать периодическое переключение всей светодиодной линейки между состояниями «включено красным цветом \verb|{1,0,0}|» и «выключено \verb|{0,0,0}|» с периодом 0{,}5~с. Алгоритм должен быть основан на использовании асинхронного таймера и логической переменной состояния.  
  \textit{Ожидаемый результат:} равномерное мигание линейки с периодом 0{,}5~с.

  \item \textbf{Эффект «бегущего огня» (один активный светодиод).}  
  Необходимо реализовать алгоритм, при котором в каждый момент времени включён только один светодиод зелёного цвета \verb|{0,1,0}|, а остальные выключены. Активный светодиод должен последовательно перемещаться от индекса \verb|0| к \verb|24| с циклическим повторением.  
  \textit{Ожидаемый результат:} визуальный эффект перемещающегося «огонька» по всей линейке.

  \item \textbf{Модель светофора на основе конечного автомата.}  
  Требуется реализовать конечный автомат с тремя состояниями: красный \verb|{1,0,0}|, жёлтый \verb|{1,1,0}| и зелёный \verb|{0,1,0}|. Состояния должны циклически сменять друг друга с заданным временным интервалом.  
  \textit{Ожидаемый результат:} периодическая и стабильная смена цветов по схеме «красный -- жёлтый -- зелёный».
\end{enumerate}

\subsection{Алгоритмические задачи}

\begin{enumerate}[resume]
  \item \textbf{Плавная пульсация яркости красного канала (PWM-модуляция).}  
  Необходимо реализовать алгоритм плавного увеличения и уменьшения яркости красного канала от значения \verb|0.0| до \verb|1.0| и обратно. Изменение яркости должно происходить малыми приращениями с использованием таймера.  
  \textit{Ожидаемый результат:} вся светодиодная линейка плавно пульсирует красным цветом.

  \item \textbf{Стробоскопический режим (белые вспышки).}  
  Требуется реализовать быстрое периодическое включение и выключение всей светодиодной линейки белым цветом \verb|{1,1,1}| с малым временным интервалом.  
  \textit{Ожидаемый результат:} частые и чёткие белые вспышки, создающие эффект стробоскопа.

  \item \textbf{Отображение чётности индекса (шахматный узор).}  
  Необходимо реализовать статический шаблон индикации, при котором светодиоды с чётными индексами светятся красным цветом \verb|{1,0,0}|, а с нечётными — синим \verb|{0,0,1}|.  
  \textit{Ожидаемый результат:} устойчивый чередующийся цветовой узор по всей линейке.

  \item \textbf{Инвертируемая полицейская мигалка.}  
  Требуется реализовать алгоритм, при котором цвета чётных и нечётных светодиодов периодически меняются местами с заданным интервалом времени.  
  \textit{Ожидаемый результат:} регулярное чередование цветовых паттернов, имитирующее работу полицейской световой сигнализации.

  \item \textbf{Генерация цветового шума.}  
  Необходимо реализовать алгоритм, который с заданной периодичностью устанавливает для каждого светодиода случайные значения компонент RGB в диапазоне от 0 до 1.  
  \textit{Ожидаемый результат:} динамично изменяющаяся цветовая картина без видимой закономерности.
\end{enumerate}

\subsection{Работа с событиями}

\begin{enumerate}[resume]
  \item \textbf{Реакция на событие посадки.}  
  Требуется реализовать обработчик событий, который при получении события \verb|Ev.COPTER_LANDED| полностью отключает светодиодную индикацию. Реализация не должна использовать блокирующие задержки.  
  \textit{Ожидаемый результат:} после посадки все светодиоды выключены.

  \item \textbf{Аварийная индикация при низком напряжении.}  
  Необходимо реализовать обработку события \verb|Ev.LOW_VOLTAGE2|, при котором все активные таймеры останавливаются, а светодиодная линейка переводится в устойчивое состояние красной индикации.  
  \textit{Ожидаемый результат:} линейка постоянно светится красным цветом, сигнализируя об аварийном режиме.

  \item \textbf{Индикатор загрузки (loading bar).}  
  Требуется реализовать алгоритм постепенного заполнения светодиодной линейки от младших индексов к старшим, при котором включённые светодиоды сохраняют своё состояние.  
  \textit{Ожидаемый результат:} визуальный эффект «заполняющейся» полосы.

  \item \textbf{Сигнал SOS по азбуке Морзе.}  
  Необходимо реализовать световую передачу сигнала «SOS» в соответствии с азбукой Морзе: три коротких импульса, три длинных импульса и снова три коротких. Длительности импульсов должны соответствовать заданным временным параметрам.  
  \textit{Ожидаемый результат:} корректно воспроизводимый световой сигнал SOS.

  \item \textbf{Двоичный счётчик секунд.}  
  Требуется реализовать алгоритм отображения текущего значения секунд в двоичном виде, используя светодиоды линейки как индикаторы отдельных битов. Обновление состояния должно происходить один раз в секунду.  
  \textit{Ожидаемый результат:} на линейке отображается двоичное представление времени в секундах.
\end{enumerate}
